\section{Provenance and distinct comparison baselines}
\label{sec:comparison}
The repository comparison is pinned to
\begin{quote}\small
\file{ef0c534c1ae1e312c61a2a9a5a3f18f5516fd42b}.
\end{quote}
The complete intake records distinguish this current repository snapshot
from the earlier immutable commits inside each inherited proof package.
An upstream source's own commit pin is not silently rewritten when a
new comparison snapshot is inspected.

Three sources of improvement must be kept separate.
\begin{enumerate}
\item The incoming CRT-anchor package supplies reuse of the encoded row
index and the range-free label proof. This round supplies the actual
weighted/quadratic growth variants and their transported exact charts.
The two-variable reduction against the Round21 principal is real, but
it does not erase the incoming anchor's different sparse/high-height
tradeoff.
\item The height-one branch keeps its original four arithmetic source
families. Monic products and chained scales change the normalization
and hence the graph, while preserving the common source zero set.
Its smallest count493 is not obtained by compiling the new direct215
source. The retained557 point is identified as inherited, not new.
\item Start-only4127 supplies the incoming shortened native coefficient
encoding. This round proves the previously missing whole-tuple transfer
to all three ten-positive charts. Its comparison is against their actual
old 295388-byte four-port tuple, not against a differently sized P66 tuple.
The incoming P66 proof in the appendix makes its own separate 6007-cell
comparison; those byte savings must not be added together.
\end{enumerate}
The consolidated sources 166,168,170 and171 are comparison inputs, not
new results of this round. Independent full coefficient recollection gives
these additional points, all on unrestricted integers:
\begin{center}\small
\begin{tabular}{@{}lrrrr@{}}\toprule
Source and chart & Variables & Degree & Terms & Height\\\midrule
166 balanced &230&4&4482&7968\\
166 compact &229&4&12006&7968\\
168 reciprocal &215&10&99914&113304\\
170 quotient carriers &217&4&19112&5184\\
171 separated radicals &217&4&25949&4576\\\bottomrule
\end{tabular}
\end{center}
These are not uniformly dominated by the new principal: for example the
balanced230 point is substantially sparser, while the217 quotient-carrier
point has lower height. The degree-ten reciprocal point is not another
quartic. The source171 separated-radical route likewise remains a distinct
comparison architecture, not an implicitly transferred scalar theorem.
Its202-variable natural-domain point has 9818 terms and height 239204765;
that is not a 202-variable unrestricted-integer equation.
The intake authenticates 890 consolidated source blobs and independently
recollects 117 incoming literal polynomials. Those finite checks do not
make the four consolidated subtrees standalone-complete: their manifests
still refer to inherited paths outside their respective subtrees. Their independent arithmetic architectures
cannot be composed with the selected fronts merely by matching a
coordinate count or a port name. Exact current source metrics and pins
are retained in the release provenance record.

The separately audited compact graphs have the following source-declared
domains and exact finite circuit metrics:
\begin{center}\small
\begin{tabular}{@{}llrrr@{}}\toprule
Source and chart & Domain & Variables & Degree & Gates\\\midrule
170 fixed residue & $\Z_{>0}$ &13&206&70\\
170 fixed residue & $\Z_{>0}$ &13&162&71\\
170 fixed residue q4 & $\Z_{>0}$ &13&154&72\\
170 rational carrier & $\Z_{>0}$ &12&438&74\\
171 separated radicals & $\Z$ &12&512&132\\\bottomrule
\end{tabular}
\end{center}
These degrees are independently attained at the authentic numerals;
operation counts include all live binary gates. The finite degree checks
are not an independent proof of each source's unbounded zero-set theorem.
In particular the 13-positive degree154 point has lower degree than the
new 10-positive degree214 point but more independent variables. Neither
metric may be omitted when comparing the two architectures.

The primary comparison is multidimensional: domain, variable count,
degree, collected support, coefficient height and live graph cost.
Changing the domain from positive to unrestricted integers requires a
proof, not a relabeling. An expanded direct quartic and a high-degree
scalar circuit have different size regimes. The report makes no
literature-priority claim or theorem that one endpoint is globally simplest.
